%==============================================================================
% Response to Reviewers -- Paper #2026MS005845
% "Controllable Probabilistic Forecasting with Stochastic Decomposition Layers"
% Journal of Advances in Modeling Earth Systems
%
% Build:  make        (latexmk -pdf response_to_reviewers.tex)
%
% STRATEGY (agreed 2026-07-27):
%   - Minimal manuscript change. No retraining; inference-only experiments OK.
%   - Where a claim outran the evidence, narrow or delete the claim.
%   - Storage (R3 C11): present BOTH archival options; user's choice.
%
% Manuscript source: manuscript/JAMES/main.tex (cloned from Overleaf).
% \tl{###} = line in that file. Red [FILL]/[ACTION] = outstanding; delete before
% submission.
%
% INTERNAL CHECKLIST (delete before submission):
%   - Key Points revision (AE + R3 marked "Accurate Key Points: No"): tl 82.
%   - Abstract tl 87 and Plain Language Summary tl 90 quote the wrong 5 MB.
%   - COI already present tl 569-570; reviewer thanks added tl 566.
%   - NOTATION FIX neither reviewer caught (do proactively): tl 346, 409 write
%     Z_beta = beta*Z and Eq. 3 interpolates Z, but the implementation scales /
%     interpolates the realized perturbation delta (scale_latents multiplies the
%     captured deltas). Proof it matters: Fig. 11 shows beta=0 -> deterministic
%     (spread -> 0), which holds only for beta*delta; scaling Z through S=WZ+b
%     leaves sigma0*R.b.M != 0 at beta=0 via the learned bias b. One-line fix:
%     define beta and interpolation as acting on the stored perturbation
%     tensors; all figures/symmetry arguments then literally correct. R3 would
%     catch this in round 2.
%   - 35 MB origin (per JS): plausibly float64 full field (31.4 MB) plus smaller
%     per-layer contributions; exact provenance not recoverable. Letter just
%     corrects the number without forensics.
%   - Gen2 status: model.yml converted -> revision_sims/model_gen2.yml via
%     `credit convert -y` (env: credit-main-casper). New sims go in
%     ensemble/revision_sims/.
%   - Only compute: one inference job (R1 minor 1 + R3 C11 verify + C5 decoder
%     timing) + scoring runs for R3 C4. No training.
%==============================================================================
\documentclass[12pt]{article}
\usepackage[margin=1in]{geometry}
\usepackage{times}
\usepackage{parskip}
\usepackage{amsmath}
\usepackage{xcolor}

% Internal markers -- delete before submission
\newcommand{\fillin}[1]{{\color{red}\textbf{[FILL: #1]}}}
\newcommand{\action}[1]{{\color{red}\textbf{[ACTION]}\ #1}}

\begin{document}

\vspace{1cm}

Dear Editor and Reviewers,

We thank the reviewers for their thorough evaluation of our manuscript. Their
feedback strengthened both the scientific content and the presentation. We
addressed every concern through revision. We are especially grateful to
Reviewer~3, whose scrutiny of our storage accounting caught a genuine error,
which we have now corrected. Below we give point-by-point responses to each
comment.

Three notes on the scope of the revision. First, Reviewer~3 asks for a
retrained baseline (Comment~1) and retrained ablations (Comment~2). Rather
than defend the claims that required them with new training runs, we removed
those claims: the paper no longer asserts a skill advantage over global noise
injection, and the design assertions that rested on unreported ``preliminary
experiments'' have been deleted or replaced with their actual rationale.
Where the underlying question can be answered at inference time from
existing forecasts, we answer it. Second, we added the multivariate
calibration diagnostics requested in Comment~4, which require only
additional scoring of forecasts already produced, not new training. Third,
correcting the storage accounting (Comment~11) turned up a problem that ran
deeper than the reviewer could have seen from the text. The corrected
presentation now gives both legitimate archival options with their true
sizes and leaves the choice to the user.

\section*{Response to the Associate Editor}

\textbf{Editor Comment:} ``While both reviewers highlight the strengths of your work,
they also raise points that need to be addressed. I therefore return the paper to you
for Major Corrections.'' (Evaluation: Accurate Key Points: No.)

\textbf{Response:} We thank the Associate Editor for handling the manuscript. We
addressed every point raised by both reviewers, and confined the revision to the
passages they identified. On the Key Points evaluation: Key Point~2 was inaccurate
on two counts identified by Reviewer~3. It claimed spread adjustment ``without
regenerating members,'' but rescaling in fact requires re-running the decoder
(Comment~5). It also rested on an incorrect storage figure (Comment~11). It now
reads: ``Ensemble members are archived as kilobyte-scale latent representations,
enabling exact regeneration and post-inference spread adjustment by re-running only
the decoder.''

\section*{Response to Reviewer 1}

We appreciate Reviewer 1's positive assessment of the work as ``a valuable result to
add to the rapidly proliferating field of ML-based weather models,'' and their
recommendation of minor revisions. Several of their questions identified genuine gaps
in our exposition. We close those gaps below.

\subsection*{Comment 1: Noise Handling During Autoregressive Rollout}

\textbf{Reviewer Comment:} ``How are these ensemble members rolled forward in time to
produce the full set of ensemble forecasts beyond 6 hours? Is the same noise injected
repeatedly for further time steps? Is a new noise tensor drawn for each 6 hour step?
Or does the deterministic model evolve those members with no additional noise
injection? This isn't clear from the text, though it is implicit that there is some
consistent rollout scheme in Table 3, but it is not clearly articulated in the text
from my reading.''

\textbf{Response:} The reviewer is right that this was never stated in the text, and
it is the most important operational detail of the method. A new, independent latent
is drawn at every 6-hour step, at each of the three injection points, so an $N$-step
forecast for one member consumes $3N$ independent draws. Noise is not held fixed along
the trajectory. Members are not evolved deterministically after the first step. The
reviewer correctly inferred this from Table~3, whose rows are forecast times, but that
table appears in Section~5.6, long after the method is introduced. We added the
explicit statement to Section~2.4 (Ensemble Generation): ``Noise is re-injected at
every autoregressive step. At each 6\,h step $t$ of the rollout, an independent latent
$\mathbf{Z}_\ell^{(t)} \sim \mathcal{N}(0,\mathbf{I})$ is drawn at each of the three
injection points $\ell$, so that an $N$-step forecast consumes $3N$ latent draws per
member. A member is identified by the complete ordered set $\{\mathbf{Z}_\ell^{(t)}\}$,
and it is this set that is archived (Section~5.4).'' We also added a forward reference
to the Table~3 caption.

\subsection*{Comment 2: Applicability Beyond U-Net-Like Architectures}

\textbf{Reviewer Comment:} ``Does this type of noise injection approach work for
architectures other than the U-Net like approach in WXFormer? It is clear that the
approach leverages the upsampling network steps to constrain the spatial structure of
the ensemble perturbations. Graph Transformer architectures don't have this same
structure \ldots{} Is the approach here only really useful on U-nets, or pure
transformers? Some sense of how these might be adapted in other contexts would
significantly improve the applicability of the method as a way to produce coherent
ensemble perturbations.''

\textbf{Response:} We added a paragraph to the Discussion (Section~6.3) on
architectural generality. The essential requirement is not the U-Net skip structure
itself, but a hierarchy of internal representations at distinct spatial resolutions.
It is the resolution of the injection site, not the latent, that sets the spatial band
a perturbation occupies. That is why Table~2 can map injection points onto
meteorological scales at all. Any encoder--decoder backbone with multi-resolution
feature maps (U-Net, U-ViT, hierarchical or Swin-style transformers) therefore
supports SDL directly. For a flat, single-resolution transformer, the scale separation
would need to come from the injection itself, for example by projecting the latent
through a band-limited basis before it enters the residual stream. For graph-based
backbones such as GraphCast or AIFS, the natural analogue is injection at each level
of the multi-mesh, with mesh refinement level playing the role of decoder resolution.
We present these two cases explicitly as reasoned expectations, not tested results.

\subsection*{Comment 3: Comparison with Loss-Based Spectral Constraints}

\textbf{Reviewer Comment:} ``There is an alternative way to constrain spatial
structure is in the loss function proposed by Lang et al.\ (arXiv:2506.10868), though
as noted in the manuscript the compact latent tensors provide an efficient way to
store ensembles that the loss-based method doesn't provide. Are there any other useful
points of comparison between these two alternative approaches, particularly around how
one might compare their frequency contributions to the forecast fields?''

\textbf{Response:} We added this comparison to Section~6.2. The distinction is this: a
spectrally-weighted loss constrains the marginal spectrum of the forecast error but
leaves the generative mechanism free, so the spectral content of any individual
perturbation is emergent and not addressable at inference. SDL instead constrains the
perturbation basis, so each latent has a known band of influence in advance. That is
precisely what makes $\beta$ a per-scale control rather than a global one. The two
approaches are complementary and could be combined. For a common quantitative point of
comparison, we suggest the kinetic energy spectrum of the perturbation, which is
computable for either method and which we already report for SDL in Figure~7. We added
Lang et al.\ (arXiv:2506.10868) to the references.

\subsection*{Comment 4: Storage Comparison with Low-Dimensional Noise Injection}

\textbf{Reviewer Comment:} ``One point in favor of the method is the ability to cache
the latent tensors and store ensembles more efficiently than the full fields and is
more useful than saving RNG seed information as shown in the manuscript. The other
alternative would be store the noise itself. For globally perturbed initial conditions
I suspect this would be more expensive than storing the compact latent tensors.
However, Alet et al.\ (arXiv:2506.10772) inject noise into the conditional norm layers
in the graph transformer, which has a much lower dimensionality --- could you comment
on how these would compare for efficient reproducibility purposes?''

\textbf{Response:} A fair challenge, and it connects directly to Reviewer~3's
Comment~11, where our storage accounting turned out to be wrong. The corrected figures
are given in full there. In brief, the storage-minimal archival recipe is three
442-dimensional latent vectors per forecast step plus the RNG seed
($3\times442\times4\,\mathrm{B}=5.3$\,kB per member per step). A self-contained
alternative, storing the realized per-layer perturbations with no seed required, is
larger, and the revised text presents both. At the kilobyte scale, we agree the size
argument no longer distinguishes SDL from conditional-normalization noise of the kind
Alet et al.\ use, which is comparably compact. We have therefore narrowed the claim. We
no longer argue that SDL's representation is smaller, only that it is structured: the
latents are indexed by injection scale, and that indexing is what supports the
per-scale $\beta$ rescaling (Section~5.5) and interpolation (Section~5.6). A flat
conditioning vector is just as cheap to archive but exposes no scale axis to
manipulate. We also softened the contrast against seed-based storage in Section~5.4,
since exact regeneration from latents still requires the seed governing the per-pixel
noise (see Reviewer~3, Comment~11). The structural argument survives, but the clean
dichotomy did not. Alet et al.\ (2025) is now cited as a comparably compact
alternative.

\subsection*{Comment 5: Reference Configuration for Figure 6 Anomalies}

\textbf{Reviewer Comment:} ``I struggle to understand how to control noise for the
different SDLs in Figure 6/Section 5.5 due to how the anomalies are computed. \ldots{}
the anomalies are computed relative to [1,1,1] (i.e. the CRPS optimal value, rather
than [0,0,0] which is just the deterministic model), so there is an asymmetry in
Figure 6 and I can't really tell if +/-3 can really be compared on this scale.''

\textbf{Response:} We agree. This is a presentation problem, and it is also the source
of the reviewer's later question in Comment~9, so one fix resolves both. We
regenerated Figure~6 with anomalies taken relative to $\beta=[0,0,0]$, the unperturbed
deterministic forecast, so positive and negative $\beta$ are directly comparable on a
common scale. This also makes the figure consistent with Section~5.7 and Figure~11,
which already use $\beta=0$ as the reference. We updated the caption and the
corresponding text in Section~5.5 to state the reference explicitly. The regenerated
figure confirms that, against the deterministic reference, the anomaly patterns are
ordered by $|\beta|$ within each column: larger magnitudes give larger-amplitude, more
spatially extensive anomalies, with consistent sign structure across the $\beta$
ladder. Positive and negative values can now be read on a common scale, as the
reviewer requested.

\subsection*{Comment 6: Generalization of Figure 6 Beyond Geopotential}

\textbf{Reviewer Comment:} ``Figure 6 shows geopotential, which is expected to be
relatively smooth in structure compared to some other fields. Do other fields respond
differently, and are the interpretations in Figure 6 generalizable?''

\textbf{Response:} A reasonable request, and one we can satisfy from existing runs. We
added the same per-layer analysis for two fields of different spectral character,
700\,hPa specific humidity and mean sea-level pressure, to the Supporting Information,
with a summary sentence in the main text. All three fields (Z500, Q700, MSLP) show the
same qualitative behavior: spatially coherent anomalies whose amplitude and extent
scale with $|\beta|$ at each injection site. This confirms that the per-layer
interpretation of Figure~6 is not specific to the smoothness of geopotential.
Figure~10 already shows five variables (T500, V500, Q700, Z500, MSLP) responding to
uniform $\beta$ scaling, so the multi-variable behavior was partly established. What
was missing was the per-layer decomposition for fields other than Z500, which the new
SI figures now supply.

\subsection*{Comment 7: Validity of Linearly Interpolated Latent States}

\textbf{Reviewer Comment:} ``The manuscript shows how one can interpolate between
latent states --- the states certainly look plausible, but I don't see any a priori
reason from the network architecture why these states should still necessarily be
valid samples drawn from the target distribution. \ldots{} Is there an explanation why
this should work, either from the architecture or CRPS metric? Or is this just an
example of how high-dimensional latent spaces work in practice?''

\textbf{Response:} The reviewer is right to press on this, and we do not over-claim in
the revision. There is no architectural guarantee, and we now say so. There is,
however, one quantifiable statement we can make. Our interpolation (Equation~3) is
linear, $\mathbf{Z}_e=(1-e)\mathbf{Z}_1+e\mathbf{Z}_2$. Because the two stored
representations are independent and zero-mean, the interpolant is zero-mean with
variance scaled by $e^2+(1-e)^2$. That factor equals one only at the endpoints and
falls to $0.5$ at $e=0.5$. The interpolant therefore has the right symmetry but is
under-dispersed, with standard deviation $\sqrt{0.5}\approx0.71$ of an independent
draw at the midpoint: the familiar norm-shrinkage effect. We now state this in
Section~5.6, along with the observation that spherical interpolation would remove the
variance deficit exactly, at the cost of no longer being a linear operation on the
archive. It also gives a concrete reading of the smooth transitions in Figure~9:
midpoint forecasts are drawn from a slightly narrower effective distribution than the
endpoints, consistent with their appearing as intermediate rather than novel states.
Separately, a valid latent draw does not guarantee that the decoded field is a valid
draw from the forecast distribution. That is an empirical property, not a guaranteed
one, and we softened the closing sentence of Section~5.6 to present it as such.

\subsection*{Comment 8: Interpretation of Figure 9 at Large versus Nominal Beta}

\textbf{Reviewer Comment:} ``Slightly confused by Fig 9 --- $\beta = \pm 3$ tend to
show broad patterns of uncertainty, while $\beta = \pm 1$ show uncertainty in specific
structural patterns that align with features in the ensemble mean. Is this just
illustrating that the models are indeed well-calibrated with $\beta = \pm 1$ and
capture the structures that produce the bulk of the uncertainty, and increasing the
ensemble spread takes them further away from that optimum? A bit more explanation
would help me understand this better.''

\textbf{Response:} The reviewer's reading is ours, and we added it to the text. At
$\beta=1$, the perturbations are the ones the afCRPS objective selected, and they
concentrate where the flow is genuinely uncertain: the baroclinic zones, jet streams,
and cyclone centers noted in Section~5.7. As $|\beta|$ grows, amplitude within each
scale band is scaled uniformly, so the flow-dependent structure is progressively
swamped by the flow-independent spatial covariance of the injection basis itself.
That produces the broad, smooth patterns. Figure~11 supports this reading
quantitatively: ensemble-mean RMSE and CRPS are both minimal near $|\beta|\approx1$
and degrade away from it. This is the same trade-off Reviewer~3 identifies in their
Comment~6, and the two passages now cross-reference each other.

\subsection*{Comment 9: Sign Symmetry of Beta in Fields versus Scores}

\textbf{Reviewer Comment:} ``Line 463 --- Yes, changing the sign of $\beta$ inverts
the perturbation in latent space, but Figure 6 shows that the anomalies are not
necessarily point-wise symmetric in the forecast field \ldots{} Are the latent
perturbations indeed spatially different between positive/negative $\beta$, but when
averaged over a full year they produce visually identical statistical scores? Or have
I missed something?''

\textbf{Response:} The reviewer has diagnosed this correctly and has not missed
anything. Both readings are true, and they are not in conflict. Because the decoder is
nonlinear, $-\beta$ does not produce the point-wise negation of the $+\beta$ field for
any individual forecast. But because the latent prior is zero-mean and symmetric, the
distribution over members is invariant under $\beta\to-\beta$, so aggregated scores
are identical up to sampling noise. The residual visual asymmetry in Figure~6 was
largely an artifact of the $\beta=[1,1,1]$ reference, now changed to $\beta=[0,0,0]$
per Comment~5. We added a sentence making the per-forecast-asymmetry versus
distributional-symmetry distinction explicit, and reworded the description of the
$\beta=\pm3$ contour ``similarity'' in Section~5.7 so that it is attributed to
ensemble statistics rather than to individual fields.

\subsection*{Minor Comment: Figure 5 and Cross-Hardware Reproducibility}

\textbf{Reviewer Comment:} ``Figure 5: I am not sure this figure adds anything ---
essentially, it shows two identical fields \ldots{} That these would be the same seems
obvious to me and the figure does not add anything. More interestingly, have you tried
reproducing the field from the stored latent tensor using different hardware or
compute configurations? GPU vs CPU, for example? This would certainly give a better
view as to how robustly reproducible this is.''

\textbf{Response:} We agree on both counts, and we ran the suggested experiment.
Figure~5 has been replaced with a cross-configuration reproducibility test:
regeneration from the stored representation on the same GPU, and on CPU, reporting
maximum absolute differences per field. This tests the property that matters
operationally, portable reproduction, rather than reproduction within a single
process, which the reviewer rightly calls obvious.

The measured outcome, over a 5-step (30\,h) rollout on an NVIDIA A100: regeneration on
the same GPU is bit-exact. The maximum absolute difference is exactly zero at every
step, both when replaying the stored perturbation tensors and when reconstructing them
from the stored latents plus RNG seed. Regeneration on CPU is not bit-exact. The
maximum absolute difference against the GPU-generated original is $2.7\times
10^{-2}$ (in the physical units of the affected prognostic field), which we attribute
to differing floating-point reduction orders between GPU (cuBLAS/cuDNN) and CPU
(MKL/oneDNN) kernel implementations. The revised text therefore states that
reproduction is exact within a fixed hardware and software configuration, and
bounded, not bitwise, across configurations. We were not able to test across different
GPU architectures, since the GPU partitions available to us at the time of revision
all carry the same device model (A100). We note this limitation in the caption.

\section*{Response to Reviewer 3}

We thank Reviewer~3 for an unusually detailed and technically precise review.
Comment~11 in particular identified a real error whose correction has improved the
paper. We are explicit below about which requests we have declined and why.

\subsection*{Comment 1 (Q1): Absence of a Global-Injection Baseline}

\textbf{Reviewer Comment:} ``A baseline against the injection scheme the paper
critiques is absent. The motivation of the work is framed against global noise
injection via conditional normalization (L100-102, 116-118, 148-149), yet that type of
injection is not implemented and tested on the WXFormer backbone. A controlled
comparison (with the same backbone, same training setup) is an important addition for
the paper's central argument.''

\textbf{Response:} The reviewer is right that the comparison is absent, and we accept
the criticism of the framing. We have not run the baseline. Instead, we removed the
claims that required it. A controlled comparison would mean fine-tuning a second model
from the same deterministic checkpoint with conditional-normalization injection, under
the same afCRPS objective and schedule, then repeating the full 2022 verification. We
do not believe the paper's contribution depends on winning that comparison. The
contribution is that SDL exposes ensemble spread as an explicit, scale-indexed,
post-hoc addressable quantity: an architectural property established by the $\beta$
and interpolation results, not a skill claim. We rewrote the three passages the
reviewer cites so that global noise injection is described as an alternative design
with different properties, not as a weaker baseline. In particular, we removed the
Abstract's assertion that the global structure ``increases training expense'' (we
retain only the interpretability distinction, which follows from where the noise
enters), and the Introduction now states plainly: ``We do not claim a skill advantage
over global conditioning. Establishing one would require an independently fine-tuned
model on the same backbone, which we have not undertaken here.''

\subsection*{Comment 2 (Q2): Ablation of Core Design Choices}

\textbf{Reviewer Comment:} ``The core design is not ablated. Decoder-only vs. full
encoder-decoder injection (L174-187), independent vs. correlated latents (L163-166),
and $\alpha$ = 0.235 (L200) are all justified only as `preliminary experiments' with
no shown sensitivity. Same for the claim that joint full-weight fine-tuning yields
superior ensemble quality over training SDL parameters alone (L231-232, 627-628),
which is stated twice with no supporting data.''

\textbf{Response:} Agreed in substance: ``preliminary experiments showed'' is not
evidence. We took three different actions, depending on what the manuscript actually
claims in each case.

First, claims we deleted. The assertion that joint full-weight fine-tuning yields
superior ensemble quality is now deleted in both places the reviewer identifies. We
have no reportable evidence for it, and we prefer deletion to softening. Joint
fine-tuning is now described simply as the configuration we used.

Second, claims that already had a stated rationale, which we sharpened rather than
removed. We would gently note that two of the three design choices are not actually
attributed to preliminary experiments in the manuscript. Decoder-only injection is
justified by computational constraints (Section~2.3). Independent latents are
justified on scale-separation grounds: synoptic/meso-$\alpha$ and meso-$\beta$
processes have different predictability characteristics, so distinct uncertainty
representations per scale are the design intent (Section~2.2). We made both
rationales unmistakable. The implementation retains switches for each configuration,
so the comparisons remain available to future work.

Third, evidence we can supply without retraining. The per-scale question, what does
each injection point contribute, is answerable at inference by single-scale
elimination: setting one scale's gain to zero removes that injection point's
contribution exactly, with fresh independent noise draws at the remaining sites. We
ran 20-member ensembles for $\beta=[0,1,1]$, $[1,0,1]$, and $[1,1,0]$ against the
$[1,1,1]$ control on the Winter Storm Izzy case. The new table reports RMSE, spread,
SSR, and CRPS at five lead times. The result is specific. Eliminating the intermediate
(meso-$\alpha$, Level~2) injection point collapses ensemble spread far more than
eliminating either of the other two. At 108\,h its SSR falls to roughly half the
control value in every field examined (e.g., Z500 0.29 vs.\ 0.44; MSLP 0.52 vs.\ 0.89;
T850 0.50 vs.\ 0.77), and it produces the worst RMSE and CRPS of the three
eliminations in every field at every lead time. Eliminating the coarsest or finest
site, by contrast, leaves both spread and skill within roughly 10\% of the control.
This is consistent with the manuscript's spectral finding that the Level-2 site has
the broadest wavenumber footprint (Figure~7c). The caption states clearly that this is
an inference-time elimination on a single case, not a retrained ablation. The two
answer related but distinct questions. For $\alpha$, see our response to Comment~3.

\subsection*{Comment 3 (Q3): Selection and Learnability of the Noise Scaling
Parameter}

\textbf{Reviewer Comment:} ``L196-197: Is $\alpha$ selection entirely trial-and-error,
or is there some intuition when choosing this parameter, given its importance
(`Appropriate $\alpha$ selection enables reasonable initial ensemble spread,
accelerating convergence of learnable SDL components')? Is there a specific reason why
it's not learnable?''

\textbf{Response:} We should simply have described what we did, and the revision now
does. The value $\alpha=0.235$ was chosen by a short sweep over candidate values,
selecting the one for which the afCRPS training loss began lowest. That is, the value
at which the initial ensemble spread was already of sensible magnitude relative to the
deterministic 6-hour error, so fine-tuning did not spend early updates correcting a
badly scaled perturbation. It was not tuned against validation skill, and we do not
claim it is optimal.

On learnability, there is a precise answer. From Equation~1, the injected perturbation
is $\alpha \cdot \mathbf{R}_{\text{Gauss}} \odot \mathbf{S}(\mathbf{Z}_\ell) \odot
\mathbf{M}$, which is linear in $\alpha$, in the learned linear map producing
$\mathbf{S}$, and in the learned channel modulation $\mathbf{M}$. Any global rescaling
of $\alpha$ is therefore absorbed exactly by $\mathbf{S}$ or $\mathbf{M}$. Making
$\alpha$ learnable adds no expressiveness, only a reparameterization of parameters the
model already learns. Its role is confined to initialization, which is exactly why
selecting it by initial loss is the appropriate criterion, and why holding it fixed
costs nothing. (We confirmed in the released checkpoint that all three noise factors
retain the value 0.235 exactly, consistent with the manuscript's statement that
$\alpha$ is non-trainable.) Section~2.3 has been rewritten accordingly.

\subsection*{Comment 4 (Q4): Multivariate and Spatial-Dependence Diagnostics}

\textbf{Reviewer Comment:} ``All calibration analysis is univariate (CRPS, SSR, rank
histograms); no multivariate or spatial-dependence diagnostic is reported, which is
the relevant test for whether the learned perturbations capture
cross-variable/spatial covariance.''

\textbf{Response:} A fair point, and it is the diagnostic most relevant to a method
whose stated appeal is structured perturbations. It requires no retraining, only
additional scoring of forecasts we had already produced, so we added two diagnostics.

First, the variogram score of order $p=0.5$ (Scheuerer and Hamill, 2015), which is
specifically sensitive to spatial dependence structure. We computed it for the SDL
ensemble over 12 initializations (first of each month of 2022, 00Z) at 24, 120, and
240\,h for Z500, T850, and Q850, on a spatially pooled grid (every sixth point, since
pairwise sums scale quadratically in the number of points). The pooling is stated in
the caption. The score grows monotonically with lead time for every field, with the
largest relative growth for Z500, consistent with its longer error-growth timescale.
One honest limitation: our archives hold IFS-ENS only as summary statistics and
regridded distribution products, not raw member fields at matching valid times. We
therefore report the SDL variogram score as a documented baseline with its lead-time
scaling, not as a head-to-head comparison, and we state this in the text.

Second, and more directly interpretable, the cross-variable spatial correlation
between $Z_{500}$ and $T_{850}$, computed per member and compared against the same
correlation in ERA5 at each of the 36 initialization--lead combinations. The SDL
members track the ERA5 coupling essentially exactly: mean member correlation 0.939
versus 0.936 in ERA5, mean absolute difference 0.004, per-combination differences
within $[-0.024,+0.016]$ with no systematic sign bias, out to 240\,h. The stochastic
perturbations do not degrade this physically meaningful cross-variable relationship
even at 10 days. This addresses the reviewer's underlying question about whether the
learned perturbations respect cross-variable covariance.

\subsection*{Comment 5 (Q5): Cost of Post-Inference Spread Adjustment}

\textbf{Reviewer Comment:} ``\,`Post-inference spread adjustment without regeneration'
(L35, 437-439, 478-480, Key Point 2): injection occurs at internal decoder layers, so
rescaling actually requires re-running the decoder, and that cost is not quantified.
The claim conflates storage savings with compute savings. The authors should clarify
this aspect.''

\textbf{Response:} The reviewer is correct, and we have withdrawn the claim as stated.
Rescaling $\beta$ requires re-running the decoder. What is avoided is the encoder
pass, the latent draw, and, the point that motivated the design, ever having to
archive the full ensemble fields. We have now measured the cost rather than leaving it
implicit. On a single NVIDIA A100 GPU (batch size 1), a full forward pass takes
228\,ms per member per step, while a decoder-only re-run from cached encoder output
takes 52\,ms. The re-run is 23\% of a full forward pass, a factor of 4.3 savings.
Extrapolated to rescaling an archived 100-member, 15-day (60-step) forecast, this is
approximately 5.2 minutes of decoder re-runs versus 22.8 minutes of full regeneration
on the same device. We corrected the phrase ``without forecast regeneration'' in all
three locations and in Key Point~2, replacing it with: ``Adjusting $\beta$ requires
re-running the decoder from the stored latents, at 23\% of the cost of a full forward
pass per member per step (measured on an A100). The encoder pass and the latent draw
are not repeated, and the full ensemble fields need never have been written to
disk.''

\subsection*{Comment 6 (Q6): Beta as a Spread--Skill Trade-Off}

\textbf{Reviewer Comment:} ``SSR remaining near unity across $|\beta|$ (L466-469) is
a consequence of both spread and ensemble-mean RMSE growing with $|\beta|$ (Fig. 10);
$\beta$ therefore trades skill for spread rather than acting as a free
calibration-preserving knob. This should be stated in the text to avoid overstating
results.''

\textbf{Response:} A correct reading of our own Figure~11, and we have adopted it. SSR
is a ratio. Holding it near unity while both numerator and denominator grow is not
calibration preservation. The underlying facts were already reported: ensemble-mean
RMSE is U-shaped with minima near $|\beta|\approx1$, and CRPS degrades symmetrically
away from it. What was missing was the interpretation the reviewer supplies, and
without it the operational recommendation of $|\beta|>1$ for expanded sampling read as
cost-free. We added: ``Because both ensemble spread and ensemble-mean RMSE grow with
$|\beta|$, the near-unity SSR across this range reflects a trade-off rather than
preserved forecast quality: CRPS degrades as $|\beta|$ departs from~1. $\beta$ should
therefore be understood as a spread--skill control, useful where deliberate
over-dispersion is wanted (e.g., tail-risk applications), and not as a free
recalibration knob.'' The operational guidance in Section~5.7 has been qualified
accordingly.

We also strengthened the empirical basis of this passage. The original Figure~11
trade-off was a single case (Winter Storm Izzy at 108\,h). We added a
multi-initialization version of the same analysis, aggregating RMSE, spread, SSR, and
CRPS as functions of uniform $\beta$ over 101 initialization dates from the 2022 test
year and 29 lead times (0--168\,h), for six variables (Z500, T850, T2m, Q700, Q500,
MSLP), over the range $\beta \in [0.25, 1.75]$. Every variable shows a U-shaped mean
RMSE and CRPS curve, minimized at or within one grid step of $\beta=1$, while SSR
increases monotonically with $\beta$ and crosses unity between roughly $\beta=1.2$
and $1.4$ depending on variable. This confirms the reviewer's reading at scale:
$\beta$ is a spread--skill trade-off, and the trained operating point sits at the
skill optimum. It replaces a single-case claim with a 101-initialization one.

\subsection*{Comment 7 (Q7): Efficiency Argument and Consistency Models}

\textbf{Reviewer Comment:} ``L82-86: The efficiency argument against diffusion rests
on its 20-50 iterative inference steps. This claim is incomplete without mentioning
single- or few-steps consistency models, which target exactly this overhead (e.g.,
Stock et al., 2025, SWIFT: An Autoregressive Consistency Model for Efficient Weather
Forecasting, NeurIPS 2025). Please compare SDL's inference cost with that of these
models as well.''

\textbf{Response:} The reviewer is right that the blanket claim is out of date, and we
revised it. Consistency and distillation approaches target exactly the overhead we
invoked, and SWIFT is a direct example. We added the citation. The revised passage
reads: ``Diffusion-based weather models have conventionally required learning full
denoising trajectories during training and tens of iterative refinement steps per
forecast at inference. Consistency and distillation approaches now reduce inference to
a few steps or one (e.g., Stock et al., 2025), so step count alone no longer separates
these methods from single-pass alternatives. The remaining distinction is whether a
separate generative model and distillation stage are required at all.'' On the
requested comparison: SDL generates a member in one forward pass of the backbone, so
against a well-distilled one-step consistency model the inference step count is
comparable rather than dramatically better, and we say so. What remains distinct is
that SDL requires no separate generative model and no distillation stage (the ensemble
capability comes from fine-tuning the existing deterministic backbone at 1.61\% of its
training cost), and that its perturbations remain addressable by scale after
inference. We report our own cost in backbone forward passes per member (one) and
quote published step counts for few-step models, marking clearly which numbers are
ours and which are quoted. We have not benchmarked SWIFT ourselves and do not present
a timing comparison against it.

\subsection*{Comment 8 (Q8): Like-for-Like Cost Accounting and Expressivity Limits}

\textbf{Reviewer Comment:} ``The claim that SDL is `substantially more efficient than
diffusion-based ensemble generation' should be clarified in two aspects. First, SDL's
reported cost is marginal: it supposes a fully pre-trained deterministic backbone, so
a like-for-like comparison to diffusion training should account for that pretraining
\ldots{} Second, efficiency is only one axis \ldots{} SDL's perturbation mechanism is
constrained to be additive and symmetric about the deterministic forecast (both R and
Z are zero-mean Gaussian, entering linearly before the residual addition). While the
downstream nonlinearity can reshape this, the architecture contains no mechanism to
produce genuinely multimodal forecast distributions because members are generated as
independent perturbations around a single deterministic trajectory rather than as
draws from a learned distribution over plausible atmospheric states. Diffusion models
instead learn an arbitrary joint distribution over forecast fields and can represent
such discrete scenario separation directly.''

\textbf{Response:} We accept both points and have revised accordingly, and we offer
one technical correction to the characterization of the mechanism that does not
change the reviewer's conclusion.

On cost accounting: we agree the marginal figure is not a like-for-like comparison
against a from-scratch generative model. The revision reports both totals and names
the framing for each: the 1.61\% fine-tuning increment, which is the operationally
relevant number for a center that already runs a deterministic model and is our
intended use case, and the pretraining-plus-fine-tuning total, which is the correct
basis for comparison against training a generative model from scratch. Appendix~A.3
already contains the forward-pass counts needed (12{,}410{,}560 baseline versus
200{,}000 fine-tuning), so this is a framing change, not new analysis.

On the mechanism: the perturbation is not linear in the two noise sources. From
Equation~1 it is $\alpha \cdot \mathbf{R}_{\text{Gauss}} \odot
\mathbf{S}(\mathbf{Z}_\ell) \odot \mathbf{M}$, where $\mathbf{R}_{\text{Gauss}}$ and
$\mathbf{Z}_\ell$ are independent Gaussians and $\mathbf{S}$ is affine in
$\mathbf{Z}_\ell$. The perturbation is therefore bilinear in the two sources: a
product of independent Gaussians, symmetric but heavy-tailed, and already
non-Gaussian before the decoder nonlinearity acts. The manuscript did not make this
structure clear enough for the reviewer to see it. That is our fault, and we clarified
the text following Equation~1.

On expressivity: the correction does not rescue the substantive point, and we concede
it without argument. SDL members are perturbations about a single deterministic
trajectory. The mechanism can represent flow-dependent, scale-structured, non-Gaussian
spread, but it has no mechanism for discrete scenario separation, such as two distinct
cyclone tracks, or cyclogenesis versus none, because there is no learned distribution
over alternative atmospheric states to draw from. A diffusion model learns the joint
distribution and can represent such bimodality directly. A new paragraph in
Section~6.1 states this limitation in the reviewer's terms, cross-referenced to the
out-of-distribution tail-risk limitation already noted there, since the two are
aspects of the same constraint. We narrowed the ``substantially more efficient''
sentence to a statement about the marginal cost of adding ensemble capability to an
existing deterministic model.

\subsection*{Comment 9 (Q9): Reconciling the Two Interpretability Claims}

\textbf{Reviewer Comment:} ``The claim `interpretable uncertainty quantification'
(L35-37) should be reconciled with the one about the method that `lacks
interpretability in terms of physical processes' (L548-549).''

\textbf{Response:} These were contradictory as written. What we meant is that the
uncertainty representation is interpretable in terms of spatial scale, since each
latent has a known band of influence, while the perturbations are not interpretable
as specific physical error sources such as convective or boundary-layer
parameterization error. The Abstract now reads ``provides a scale-interpretable
representation of forecast uncertainty,'' and we added a clause distinguishing the two
senses to the Discussion passage.

\subsection*{Comment 10 (Q10): Evaluation Ensemble Size}

\textbf{Reviewer Comment:} ``The ensemble size for evaluation is not stated in the
text (x-axis in Figs. 2 and 3 imply about 100 members for SDL vs. 50 for IFS-ENS;
training used M = 10, L235). This should be clarified.''

\textbf{Response:} The reviewer's inference is correct: 100 members for SDL-WXFormer,
50 for IFS-ENS, and $M=10$ during training. We should have stated it. The sizes are
now given in Section~5.1 and in the captions of Figures~2 and~3. We also address the
fairness question this implies. We added a sentence noting which diagnostics are
ensemble-size sensitive (spread, SSR, rank histograms) and which are approximately not
(the almost-fair CRPS, which corrects for finite ensemble size via Equation~2), and we
report SDL scores subsampled to 50 members alongside the 100-member scores, so the
IFS-ENS comparison is not confounded by size. The subsampling comparison confirms the
insensitivity directly. Over 12 initializations (the first of each month of 2022,
00Z) at 24, 120, and 240\,h lead times, reducing the Z500 ensemble from 100 to 50
members changes area-weighted ensemble-mean RMSE by less than 0.6\%, spread by less
than 0.4\%, SSR by less than 0.9\%, and CRPS by less than 0.7\% at every lead time
tested.

\subsection*{Comment 11 (Q11): Archived Representation and Storage Arithmetic}

\textbf{Reviewer Comment:} ``The authors claim that they store `latent tensors'
(L188-191), whereas at L351-354 they say they store the perturbation fields. Also,
the `$\sim$35 MB per member' full-field baseline (line 354) does not reconcile with
71x192x288 float32 $\approx$ 15.7 MB. Please clarify what is archived and recompute
the comparison.''

\textbf{Response:} The reviewer is right on every count, and the problem runs deeper
than the text revealed. We are grateful for the scrutiny. Correcting it has
strengthened the paper's claim considerably.

The inconsistency was real. Section~2.3 describes storing latent tensors, while
Section~5.4 states that the computed noise perturbations are archived instead. These
are different objects, and both of our numbers were wrong. The published
$\sim$5\,MB figure applied the output-variable channel count (71) to the decoder
grids, which matches neither object. The realized perturbations live on the decoder
feature maps (1024, 512, and 256 channels on the internally padded grid), giving
\[
  1024{\times}36{\times}48 + 512{\times}72{\times}96 + 256{\times}144{\times}192
  = 1.24\times10^{7}\ \text{floats} = 49.5\ \mathrm{MB}
\]
per member per step: roughly ten times larger than reported, and, as the reviewer's
corrected baseline shows, larger than the forecast fields themselves
($71{\times}192{\times}288$ float32 $=15.7$\,MB, not $\sim$35\,MB). The comparison as
published pointed the wrong way.

The revision resolves this by presenting the two legitimate archival options
explicitly and leaving the choice to the user, since which is preferable depends on
the use case:

\emph{Option A: raw latents plus RNG seed (storage-minimal).} Each latent is
442-dimensional, one per injection point per forecast step:
$3\times442\times4\,\mathrm{B}=5.3$\,kB per member per step, a factor of roughly
$3\times10^{3}$ below the full fields. The caveat, now stated: the latents alone do
not determine the member, because the per-pixel noise $\mathbf{R}_{\text{Gauss}}$ is
drawn inside each layer. Exact regeneration under Option~A requires the latents
together with the RNG seed, a handful of integers, negligible in size, but part of
the recipe.

\emph{Option B: realized perturbation tensors (self-contained).} Archiving the
per-layer perturbations after $\mathbf{R}_{\text{Gauss}}$ is drawn removes the seed
dependence entirely, at a cost of 49.5\,MB per member per step. This is the
representation our analysis tooling used, so the paper's $\beta$-rescaling and
interpolation results were in fact produced under Option~B. Under Option~A, the
identical manipulations are available after reconstructing the perturbations from
the latents and seed.

Both options require the decoder re-run quantified under Comment~5 for any post-hoc
manipulation. We verified both directly before making the claim. Over a 5-step
(30\,h) rollout on an NVIDIA A100, replay from stored perturbation tensors (Option~B)
and replay from stored latents plus RNG seed with the per-pixel noise re-drawn
(Option~A) each reproduce the original ensemble member bit-exactly: maximum absolute
difference of exactly zero at every step for both paths. Sections~2.3 and~5.4 have
been made consistent, the arithmetic is shown explicitly, the corrected figures
replace the 5\,MB value in the Abstract, Plain Language Summary, and Key Point~2, and
the baseline is corrected to 15.7\,MB.

\subsection*{Comment 12 (Q12): Test Period Consistency}

\textbf{Reviewer Comment:} ``L221-222 state testing on 2020-2022; L247 and L687-688
state 2022. This should be reconciled.''

\textbf{Response:} Both statements are true, but the distinction was never drawn.
2020--2022 is the held-out test split following Schreck et al.\ (2025), while all
results reported here verify against 2022 only, on 00Z initializations. Section~3 now
reads ``testing (2020--2022), of which the 2022 00Z initializations are used for all
verification reported here.'' 2022 is used consistently wherever results are
reported.

\subsection*{Comment 13 (Q13): Symbol Collision for the Noise Scaling Parameter}

\textbf{Reviewer Comment:} ``Since $\alpha$ is already used in Eq. 2 as afCRPS
hyperparameter, the authors may want to use using a different noise-scaling scalar for
SDL in eq. 1.''

\textbf{Response:} Agreed. The collision is genuine and confusing, since
$\alpha=0.235$ (Equation~1) and $\alpha=0.95$ (Equation~2) appear within two pages of
each other. We renamed the SDL noise-scaling scalar in Equation~1 to $\sigma_0$ and
updated it throughout the text, Table~1, and captions. The afCRPS hyperparameter
retains $\alpha$.

\subsection*{Comment 14 (Q14): Training Data Specification}

\textbf{Reviewer Comment:} ``Training dataset: although the dataset is the same as
that used in Schreck et al., 2025 and Sha et al., 2025b, it is still worth specifying
the temporal resolution and spatial domain.''

\textbf{Response:} We extended Section~3 to specify the 6-hourly temporal resolution,
the global domain, the horizontal grid used for training, and the composition of the
74 input variables across upper-air levels, surface fields, static fields, and
forcing. The 16 selected model levels and the 1979--2017 / 2018--2019 / 2020--2022
split were already stated and are retained.

\subsection*{Comment 15 (Q15): Grid Specification Consistency}

\textbf{Reviewer Comment:} ``Section 3 describes the training data as ERA5 coarsened
to 1 degree (360x181), whereas Appendix B1 states training was performed on a
regridded 192x288 grid, which is also the evaluation grid. Please reconcile and
clarify: was the model trained natively on 192x288 (and 360x181 referenced only as
the source archive), and is 192x288 used end-to-end for fine-tuning and
verification?''

\textbf{Response:} The reviewer's proposed reading is correct, and the inconsistency
was real. The 192$\times$288 grid (approximately 0.94$^\circ$ latitude $\times$
1.25$^\circ$ longitude, the CESM finite-volume grid) is used end-to-end for training,
fine-tuning, inference, and verification. We revised Section~3 to give the data path
once and in order: the ERA5 native 0.25$^\circ$ archive, regridded using conservative
interpolation (following the approach of WeatherBench2, as described in the cited
data-processing references, Sha et al., 2025a,b), to the single 192$\times$288 grid
used throughout. We removed the 360$\times$181 figure, or relabeled it explicitly as
an intermediate archive resolution rather than the training grid, cross-referenced to
Appendix~B.1.

\subsection*{Comment 16 (Q16): Introducing the Role of Beta at First Appearance}

\textbf{Reviewer Comment:} ``L364-367: The authors may want to specify the role of
$\beta$ values when they are introduced at L185-187, clarifying that they scale the
noise and explaining the purpose of using them. This would improve clarity.''

\textbf{Response:} Agreed. The symbols $\beta_1,\beta_2,\beta_3$ first appear in
Table~2 and Section~2.3 as labels for injection points, but their function is not
defined until Section~5.5. We added at first appearance: ``The parameters $\beta_1$,
$\beta_2$ and $\beta_3$ are per-scale gains applied at inference time to the
stochastic perturbation injected at each point. They are fixed at unity throughout
training and varied only afterwards, which is what allows ensemble spread to be
adjusted by spatial scale on an already-generated forecast (Section~5.5).'' (Wording
chosen to match the implementation, which scales the injected perturbation. See the
notation note in the preamble checklist.)

\subsection*{Comment 17 (Q17): Introduction Definition Too Terse}

\textbf{Reviewer Comment:} ``L119-121: That definition is rather too concise and
cryptic for the introduction section where the reader is still not fully aware of the
architecture.''

\textbf{Response:} Agreed. The sentence introducing SDL as layers that ``separate
latent-driven modulation from per-pixel stochasticity'' presumed the reader already
knew the decomposition being referred to. We expanded it into three sentences of
plain description that do not presume knowledge of the WXFormer decoder: what is
injected, where it enters the network, and why the separation matters for scale
control. The formal definition is left to Section~2.2.

\vspace{1em}

We thank both reviewers again for their careful reading. It has materially improved
the manuscript.

\end{document}
